\documentclass[12pt]{article}
\usepackage[utf8]{inputenc}
\usepackage[T1]{fontenc}
\usepackage{amsmath,amssymb}
\usepackage{geometry}
\geometry{margin=2.5cm}

\begin{document}

\noindent\textbf{(5)} Transportul unei legi de compoziție printr-o bijecție

\[
(A,\varphi) \qquad A \xrightarrow{f} A'
\]

\noindent Dacă $x',y'\in A' \Rightarrow \exists!\, x,y\in A$ a.î.\ $x'=f(x)$ şi $y'=f(y)$

\[
\varphi' : A'\times A' \to A', \qquad \varphi'(x',y') = f\big(\varphi(x,y)\big) \in A'
\]

\noindent [notație: $\varphi(x,y)=xy$] \quad $x'y' \overset{\text{def}}{=} f(xy)$\footnote{această linie e foarte compactată în original; se pare că se introduce convenția multiplicativă $\varphi(x,y)=xy$, astfel încât definiția lui $\varphi'$ devine $x'y'=f(xy)$.}

\[
f : A\to A' \qquad f(xy) = f(x)f(y) \quad (\forall)\,x,y\in A.
\]
\[
f^{-1}(x'y') = f^{-1}(x')f^{-1}(y') \quad (\forall)\,x',y'\in A'.
\]

\bigskip
\noindent\underline{Teoremă:} Fie $(A,\varphi)$, $\varphi:A\times A\to A$.

\noindent $f:A\to A'$ şi $\varphi':A'\times A'\to A'$.

\begin{enumerate}
\item[1)] $\varphi$ = asoc.\ (comut.) $\Leftrightarrow$ $\varphi'$ e asoc.\ (comut.).
\item[2)] $e\in A$ elem.\ neutru pt.\ $\varphi$ $\Leftrightarrow$ $f(e)$ el.\ neutru pt.\ $\varphi'$.\footnote{în original apare ,,pt.\ $f$'' / ,,pt.\ $f$' '' -- corectat aici la $\varphi$/$\varphi'$ pentru consecvență cu punctul 1) şi cu restul teoremei.}
\item[3)] $x\in A$ simetrizabil în raport cu $\varphi$ $\Leftrightarrow$ $f(x)$ simetrizabil în raport cu $\varphi'$.
\end{enumerate}

\noindent $x',y'\in A'$.

\[
x'y' = f(x)f(y) = f(xy) = f(yx) = f(y)f(x).
\]

\bigskip
\begin{center}
\fbox{\Large\textbf{GRUPURI}}
\end{center}

\noindent\underline{\textbf{Subgrupuri. T.\ lui Lagrange}}

\bigskip
\noindent\underline{Def:} Fie $G$ un grup. O submulțime $H$ nevidă a lui $G$ s.n.\ \textbf{subgrup} al grupului $G$ dacă:

\begin{enumerate}
\item[(1)] $(\forall)\,x,y\in H \Rightarrow xy\in H$
\item[(2)] $(\forall)\,x\in H \Rightarrow x^{-1}\in H$
\end{enumerate}

\bigskip
\noindent\underline{Pb:} Dat un grup $G$ să determinăm toate congruențele pe $G$.

\noindent Rel.\ de ech.\ pe $G$ compatibilă cu operația grupului:

\[
xRy \Rightarrow zxRzy.
\]

\end{document}
